% SPDX-License-Identifier: Apache-2.0
\section{What Each Language Is}
\label{sec:two-languages}

\subsection{LHDL: dataflow, and time as an implementation detail}

LHDL starts from twelve theses. The eighth is the one that shapes the
language: at the design level there should be no notion of combinational
versus sequential logic. The designer writes dataflow and never counts
cycles, and the compiler decides where registers go.

The mechanism is the \term{tag}. A tag names a synchronisation domain.
Operations sharing a tag are aligned by the compiler through longest path
analysis, which inserts bridge registers into the shorter paths. Adding
\code{with handshake} makes the domain elastic: the compiler replaces
flip-flops with skid buffers and wires ready and valid across the whole
domain. Adding \code{capacity=32} makes the compiler generate a scoreboard
that tracks up to 32 transactions in flight and asserts backpressure when
full.

The organising idea is the \term{three facets}. The design facet states
logic and contracts and is timeless. The implementation facet maps logic
onto fabric and handles retiming. The configuration facet binds
implementations into sockets, sets generic parameters, maps tags onto
physical clocks, and links foreign code. Changing the target fabric edits
the configuration facet and leaves the design facet untouched.

Reuse runs through \term{interfaces} with \term{modports}. An interface
declares fields, and a modport names a role and gives each field a direction
for that role. A module, a pipeline or a state machine can implement an
interface.

Verification is delegated. LHDL adds no non-synthesisable subset. The
configuration facet instead binds an interface to a foreign implementation
in Go, Rust or C++, and the compiler generates a gRPC bridge and bit-exact
software structs.

\subsection{TxHDL: transactions, and time written down}

TxHDL starts elsewhere. Communication between modules happens through
transactions over buses, never through raw signals. Inside a module one
sequential flow is easy to reason about, while modules run in parallel.
Variables become registers based on how they are used.

The mechanism is \code{await}. A sequence blocks until a transaction
arrives, until a condition holds, or for a stated number of cycles.
\code{select} waits on several events and takes the first. A timeout bounds
the wait. A protocol that takes 40 cycles reads as 40 lines of straight
code rather than as a state enumeration.

Communication runs through \term{buses} made of \term{channels}, each with
implicit valid and ready handshaking. A channel marked \code{tagged by id}
matches responses to requests out of order, and a \code{then} clause
registers a callback that fires when the matching response arrives.

Typing is \term{structural}. Two transactions with the same fields are
compatible whatever their names, and a transaction with extra fields can be
sent where fewer are expected.

Verification is native. Assertions, assumptions, cover points and
properties are in the language, and properties use temporal operators.

\subsection{The asymmetry}

Four of the seven concerns in Figure~\ref{fig:asymmetry} are worked out on
both sides. The three that are not decide the merge, and each is thin on a
different side.

LHDL has no way to write a protocol body. Its state machine is a hand
written state enumeration, which is the form \code{await} exists to remove.
TxHDL has no way to assemble a system. Its top level offers instances,
connections and an address-decoding bus, and there is no interface, no late
binding, no configuration facet, no way to wrap legacy Verilog and no way to
swap an implementation.

\begin{figure*}[t]
\centering
\begin{tikzpicture}[
    node distance=4mm and 10mm,
    row/.style={font=\footnotesize\bfseries,align=right,text width=26mm},
    cell/.style={box,text width=52mm,align=left},
    % Two channels, not one. An IEEE article gets printed and photocopied,
    % and a 3 per cent grey difference does not survive that. A worked-out
    % mechanism is shaded and solid; one that is only named is unshaded and
    % dashed, so the distinction reads even with the fill gone.
    thin/.style={cell,fill=white,draw=framegray,dashed},
    thick/.style={cell,fill=cellshade,draw=black},
  ]

  \node[font=\footnotesize\bfseries] (hl) at (0,0) {LHDL};
  \node[font=\footnotesize\bfseries,right=62mm of hl] (ht) {TxHDL};
  \node[row,left=10mm of hl] (h0) {};

  \node[row] (r1) at ($(h0)+(0,-9mm)$) {system\\assembly};
  \node[thick,right=10mm of r1] (l1)
    {three facets; \code{config} and \code{bind}; domain to clock mapping};
  \node[thin,right=10mm of l1] (t1)
    {\code{instance}, \code{connect}; address-decode bus};

  \node[row,below=of r1] (r2) {contracts};
  \node[thick,right=10mm of r2] (l2)
    {\code{interface} plus \code{modport}; \code{implements}};
  \node[thick,right=10mm of l2] (t2)
    {structural typing on transaction fields};

  \node[row,below=of r2] (r3) {inter-module\\communication};
  \node[thick,right=10mm of r3] (l3)
    {tags, chaining, fork and join; elastic buffers; scoreboards};
  \node[thick,right=10mm of l3] (t3)
    {\code{bus} plus \code{channel}; valid and ready; callbacks};

  \node[row,below=of r3] (r4) {intra-module\\behaviour};
  \node[thin,right=10mm of r4] (l4)
    {\code{pipeline} plus \code{stage}; \code{state} blocks};
  \node[thick,right=10mm of l4] (t4)
    {\code{seq}, \code{await}, \code{select}; \code{comb}, \code{wire},
     \code{always}};

  \node[row,below=of r4] (r5) {types};
  \node[thick,right=10mm of r5] (l5)
    {\code{u32}, \code{uint<W>}};
  \node[thick,right=10mm of l5] (t5)
    {\code{u<N>}, \code{i<N>}, \code{bits<N>}, arrays, enums};

  \node[row,below=of r5] (r6) {verification};
  \node[thick,right=10mm of r6] (l6)
    {foreign function interface; gRPC co-simulation};
  \node[thick,right=10mm of l6] (t6)
    {\code{assert}, \code{assume}, \code{cover}, \code{property}};

  \node[row,below=of r6] (r7) {grammar};
  \node[thick,right=10mm of r7] (l7) {full EBNF, LL(1) claimed};
  \node[thin,right=10mm of l7] (t7) {none};

\end{tikzpicture}
\caption{Where each language is detailed. A shaded cell is a worked-out
mechanism; an unshaded one is named and not specified. Four rows are worked
out on both sides. The three that are not, system assembly, intra-module
behaviour and the grammar, are the ones that make the merge possible: each
is thin on one side and thick on the other, and the thin side differs.}
\label{fig:asymmetry}
\end{figure*}
